\documentclass[11pt,a4paper]{article}
\usepackage[T1]{fontenc}
\usepackage[utf8]{inputenc}
\usepackage{lmodern}
\usepackage[margin=25mm,headheight=14pt]{geometry}
\usepackage{amsmath,amssymb,amsthm,mathtools}
\usepackage{booktabs,array,tabularx,longtable}
\usepackage{microtype,enumitem,xcolor,graphicx}
\usepackage{xurl}
\usepackage{fancyhdr}
\usepackage[colorlinks=true,linkcolor=blue!45!black,urlcolor=blue!50!black,citecolor=blue!45!black]{hyperref}
\hypersetup{pdftitle={Structural reductions for explicit Diophantine consistency equations},pdfauthor={Prepared for Vladimir Reshetnikov with AI assistance},pdfsubject={14 positive variables of degree 30; 41 integer variables of degree 4; a smaller nine-variable packing}}
\newcommand{\Z}{\mathbb Z}
\newcommand{\N}{\mathbb N}
\newcommand{\Con}{\operatorname{Con}}
\newcommand{\Pos}{\operatorname{Pos}}
\newcommand{\ZFC}{\mathsf{ZFC}}
\newcommand{\ZFCRef}{\operatorname{Ref}_{\ZFC}}
\newcommand{\pop}{\operatorname{pop}}
\newtheorem{theorem}{Theorem}[section]
\newtheorem{lemma}[theorem]{Lemma}
\newtheorem{proposition}[theorem]{Proposition}
\newtheorem{corollary}[theorem]{Corollary}
\theoremstyle{definition}
\newtheorem{definition}[theorem]{Definition}
\theoremstyle{remark}
\newtheorem{remark}[theorem]{Remark}
\setlist{itemsep=3pt,topsep=5pt}
\setlength{\parindent}{0pt}
\setlength{\parskip}{5pt plus 1pt minus 1pt}
\allowdisplaybreaks[2]
\emergencystretch=2em
\pagestyle{fancy}
\fancyhf{}
\fancyhead[L]{\small Explicit Diophantine consistency equations}
\fancyhead[R]{\small Research contribution, 23 October 2026}
\fancyfoot[C]{\thepage}

\title{\vspace{-12mm}\textbf{Structural reductions for explicit\\Diophantine consistency equations}\\[5mm]
\large An adaptive Pell modulus, delayed sign recovery,\\and a smaller certified nine-variable packing}
\author{Prepared for Vladimir Reshetnikov and the ProveIt repository\\
\small AI-assisted mathematical research; internal review, not external peer review}
\date{23 October 2026}

\begin{document}
\maketitle
\thispagestyle{empty}
\begin{abstract}
We give a fixed polynomial of exact degree 30 in fourteen strictly positive
integer variables, with an 83-operation division-free evaluator, and a fixed
quartic in 41 unrestricted integer variables.  For each, ordinary mathematical
proofs composed with the explicitly identified native compiler and machine
calibration establish that ZFC proves its unsolvability equivalent to
syntactic consistency of ZFC.  The first construction combines an inherited
integral-root coordinate with a new adaptive index modulus $U=Y^2$; the
divisibility it enforces is proved adequate before the deleted quotient is
recovered.  The second delays the sign of the radix until a global power of
two is recovered, then uses a scale factor of four to replace two independent
positivity representations by one product representation.  Exact elimination
gives companion sextic and octic equations in 36 and 31 unrestricted integers.
A separate finite packing certificate improves the inherited nine-positive
construction to exact degree $1{,}781{,}068$ in 470 operations.

The archive supplies every polynomial, fixed integer coefficient, generator,
source pin, proof, and finite verification receipt.  These results improve
specific measured tradeoffs; no global minimum, proof of consistency, or new
end-to-end proof-assistant certificate is claimed.
\end{abstract}

\vspace{3mm}
\begin{center}
\begin{tabular}{lrrrr}
\toprule
Construction & Unknowns & Domain & Exact degree & Operations\\
\midrule
Adaptive modulus $P_{14}$ & 14 & $\Z_{>0}$ & 30 & 83\\
Signed quartic $Q_{41}$ & 41 & $\Z$ & 4 & ---\\
Eliminated companion $Q_{36}$ & 36 & $\Z$ & 6 & ---\\
Eliminated companion $Q_{31}$ & 31 & $\Z$ & 8 & ---\\
Certified packing $F_9$ & 9 & $\Z_{>0}$ & $1{,}781{,}068$ & 470\\
\bottomrule
\end{tabular}
\end{center}
The quartic and its companions are supplied as sums of 26, 21, and 16
squared residuals, respectively.  No optimized gate count is asserted for
those three presentations.

\clearpage
\tableofcontents
\clearpage

\section{The result, its scope, and the comparison}
\label{sec:scope}

\subsection{What is being proved}

Fix the repository's recursive predicate for finite syntactic ZFC
refutations, and write $\ZFCRef(p)$ for that predicate.  In this paper
\[
\Con(\ZFC)\quad\text{means}\quad
\neg\exists p\in\N\;\ZFCRef(p).
\]
This fixes a proof-theoretic statement, rather than an informal assertion of
truth of the axioms.  Every coefficient in the equations below is a fixed
integer defined by finite data included in the archive.  None depends on
whether a refutation exists.

\begin{theorem}[Consistency equations]
\label{thm:main}
Let $P_{14},Q_{41},Q_{36},Q_{31},F_9$ be the literal polynomials defined in
the accompanying components and described below, with their complete bound
coefficient tuples.  The ordinary mathematical reductions in this paper and
the identified inherited compiler and calibration results prove in ZFC that
\begin{align*}
\neg\exists\mathbf x\in\Z_{>0}^{14}\;P_{14}(\mathbf x)=0
&\ \longleftrightarrow\ \Con(\ZFC),\\
\neg\exists\mathbf x\in\Z^{41}\;Q_{41}(\mathbf x)=0
&\ \longleftrightarrow\ \Con(\ZFC),\\
\neg\exists\mathbf x\in\Z^{36}\;Q_{36}(\mathbf x)=0
&\ \longleftrightarrow\ \Con(\ZFC),\\
\neg\exists\mathbf x\in\Z^{31}\;Q_{31}(\mathbf x)=0
&\ \longleftrightarrow\ \Con(\ZFC),\\
\neg\exists\mathbf x\in\Z_{>0}^{9}\;F_9(\mathbf x)=0
&\ \longleftrightarrow\ \Con(\ZFC).
\end{align*}
The respective exact total degrees are $30,4,6,8,1{,}781{,}068$.
\end{theorem}

The positive statements proved first are existence equivalences:
an equation has a zero in its specified domain if and only if a finite
ZFC refutation exists.  Taking negations gives Theorem~\ref{thm:main}.
This does not prove either consistency or inconsistency.  It also does not
conflict with the second incompleteness theorem: an equivalence to
consistency is different from a proof of either equivalent assertion.

``Provable in ZFC'' here has the ordinary mathematical meaning: all finite
syntax, arithmetic, induction, and compiler arguments used in the proof can
be carried out in ZFC.  We do not supply a checked object-language ZFC
derivation, or a new closed Lean/Rocq theorem for these literal equations.
The repository already distinguishes those levels of evidence
\cite{active}; the distinction is maintained throughout this delivery.

\subsection{Pinned and comparable baselines}

The review is pinned to ProveIt commit
\begin{center}\small\texttt{beb0192cb8dc36f78c369c080f2adfbc1887a486}.\end{center}
The general collection README still describes earlier waves at that pin.
We compared against the later component files as well, including packages
55, 57, 58, and 59.  The relevant comparisons are as follows
\cite{source55,source57,source58,source59}.

\begin{center}
\small
\begin{tabularx}{\textwidth}{@{}l*{3}{>{\raggedright\arraybackslash}X}@{}}
\toprule
Criterion & Inspected predecessor & This contribution & Change\\
\midrule
Fourteen positive variables
& Degree 52, 79 operations (55/59)
& Degree 30, 83 operations
& $42.31\%$ lower degree, four extra gates\\
Positive degree 30
& Fifteen variables, 79 operations (58)
& Fourteen variables, 83 operations
& One fewer variable, four extra gates\\
Unrestricted integer quartic
& 45 variables, 28 quadratic rows (57)
& 41 variables, 26 quadratic rows
& Four fewer variables, two fewer rows\\
Nine positive variables
& Degree $1{,}910{,}404$, 473 operations (59)
& Degree $1{,}781{,}068$, 470 operations
& $6.77\%$ lower degree, three fewer gates\\
\bottomrule
\end{tabularx}
\end{center}

The sextic and octic companions improve the named historical 42-variable
sextic and 36-variable octic presentations in packages 27/29, as retained
and identified in the later collection.  We do not claim an exhaustive
comparison of all subsequent intermediate low-degree variants.  The new
fourteen-variable form does not dominate the fifteen-variable degree-24
form of package 58 or the fifteen-variable 66-operation form of package 51.
Likewise a quartic in 41 integers does not improve the separate
minimum-variable high-degree integer constructions.

There is no specified total order on ``simplest.''  Variable count, witness
domain, degree, number of arithmetic gates, and sizes of fixed numerals are
different measures.  We therefore give the explicit tradeoffs rather than
claim a global optimum.  The main new mathematical content is the adaptive
modulus and the two sign reductions.  The nine-variable improvement is a
smaller certified finite layout.

\subsection{Conventions for explicitness and cost}

A register definition is an abbreviation for a polynomial expression,
unless the variable is explicitly listed among the quantified coordinates.
In particular a displayed chain of definitions never silently adds
existential variables.  Expanding the abbreviations gives a member of
$\Z[x_1,\ldots,x_n]$.  There are no variable exponents, divisions,
square roots, bit operations, or run-time sequence evaluations in the
delivered polynomials.  Such operations may occur in their proofs.

The gate metric charges every live binary addition, subtraction, and
multiplication, including multiplication by a fixed numeral.  Squaring
costs one multiplication.  Reuse and fixed leaves are free, and
constant-only arithmetic is folded.  This is an arithmetic-circuit metric,
not bit complexity or a proof of circuit optimality.  Exact degree is
measured after polynomial cancellation, with each unknown of degree one
and every fixed numeral of degree zero.

The fixed coefficients are intentionally given as exact binary payloads
and finite generators instead of several million decimal digits.  If a
decompressed file has bytes $b_0,\ldots,b_{m-1}$, its integer is
$\sum_{j=0}^{m-1}b_j256^j$.  The compression is lossless; the byte strings,
their hashes, and the generators all accompany the paper.  This is an
exact definition of each numeral, with no oracle or unbounded search.

\subsection{The inherited semantic endpoint}

All three families retain the native coefficient/compiler route of the
inspected repository.  The new fourteen-positive and signed low-degree
reductions act directly on its Pell arithmetic; they do not invoke a
universal nine-variable polynomial.  The additional nine-variable form
does retain the inherited Matiyasevich--Jones finalizer.

The fixed endpoint is the calibrated blank-tape machine K215.  The
repository identifies the theorem
\begin{center}
\small\path{BusyBeaver.ZFC.KEndpoint215.blank_halts_iff_zfc_refutation}
\end{center}
and the syntactic endpoint \texttt{BProv ZFCax\_s [] Form.fBot}
\cite{active}.  The native-window construction maps its finite accepting
computations to the fixed arithmetic source and conversely.  Each full
coefficient tuple has its own adapter and canonical completeness witnesses.
The transformations below preserve this endpoint.  Although a route
independent of this machine was considered, no smaller fully specified and
verified alternative endpoint was obtained; no such result is claimed.

\section{Fourteen positive unknowns at degree thirty}
\label{sec:sc-fourteen}

The construction in this section combines a free Pell ordinate with a new
adaptive quotient modulus.  It gives a fixed polynomial with fourteen
strictly positive integer unknowns, exact total degree $30$, and a
division-free evaluator using $83$ binary arithmetic operations.  The
principal improvement is the modulus $U=Y^2$: it has polynomial degree two,
but a preceding integrality argument forces it to be large enough for the
index recovery.  The proof establishes that size before invoking the
strong Pell-rank theorem.

\begin{theorem}\label{thm:sc-main}
Let $P$ be the polynomial defined in
\eqref{eq:sc-registers}--\eqref{eq:sc-polynomial}, with the complete fixed
coefficient tuple specified below.  Then
\[
 \mathrm{ZFC}\vdash
 \left(\exists\boldsymbol v\in\Z_{>0}^{14}\ P(\boldsymbol v)=0\right)
 \ \longleftrightarrow\
 \left(\text{there is a finite syntactic ZFC refutation}\right).
\]
Consequently
\[
 \mathrm{ZFC}\vdash
 \left(\forall\boldsymbol v\in\Z_{>0}^{14}\ P(\boldsymbol v)\ne0\right)
 \ \longleftrightarrow\ \Con(\mathrm{ZFC}).
\]
The exact degree is $30$.  The supplied arithmetic circuit has $44$
multiplications, $23$ additions, and $16$ subtractions.
\end{theorem}

Here and below, the semantic conclusion composes the new elementary
reduction with the inherited native compiler and the literal
K215/refutation calibration.  The full ordinary mathematical argument and
its source dependencies are supplied; an additional end-to-end proof
assistant certificate is not claimed.  No consistency assumption enters
either direction.

All artifact paths in this section are relative to \path{positive14/}.

\subsection{The fixed polynomial and its coefficient binding}

The fourteen supplied coordinates are
\[
 (J,F,\alpha,z,f,i,T,Y,w,\rho,\eta,\zeta,e,\gamma)
       \in\Z_{>0}^{14}.
\]
All further letters in the definition are polynomial abbreviations, not
additional quantified variables.  Use the \emph{entire} compact native
coefficient tuple from package 44.  Its four ports are
\[
\begin{gathered}
 B=2^{590069},\qquad
 \beta=B-1,\qquad K_0=D_C+BD_R,\\
 M_C=M_{C,\mathrm{native}},\qquad
 M_F=M_{F,\mathrm{native}}+B-1.
\end{gathered}
\]
In particular, $M_F$ is the already shifted source port.  The exact
integers are the unsigned little-endian byte streams obtained by
decompressing the following four delivered files.  Their hashes, finite
coefficient recipes, and binding are recorded in
\path{native/COEFFICIENT_BINDING.json}.
\begin{center}
\begin{tabular}{lll}
\hline
Port & Exact payload under \path{native/coefficients/} & Raw bytes\\
\hline
$\beta$ & \path{beta.bin.gz} & $73\,759$\\
$K_0$ & \path{K.bin.gz} & $74\,111$\\
$M_C$ & \path{MC.bin.gz} & $73\,759$\\
$M_F$ & \path{MF_source.bin.gz} & $73\,759$\\
\hline
\end{tabular}
\end{center}
The coefficient recipe uses radix $2^{67}$ and length $8807$, with the
ordered literal list of $943$ Wang tiles and the $430$ literal K215
transition rows included in \path{native/}.  The independent generator and
auditor are included there too.  Thus this is one fixed four-tuple, rather
than four freely adjustable parameters.  In particular it satisfies
\[
 0<M_C<B-1,\qquad B-1<M_F<2(B-1),
 \quad M_C\equiv2\pmod4,
 \quad M_{F,\mathrm{native}}\equiv4\pmod8,
\]
and the native mask populations sum to $590069$.  Its complete native
interface, including canonical completeness, is proved in
\path{native/PROOF.md}.  The new construction binds all four ports
together to this interface.

Define
\begin{equation}\label{eq:sc-registers}
\begin{aligned}
 q&=\beta J+1,                 & C&=q-F-\alpha,\\
 X&=wq,                       & E&=XY,\\
 k&=\eta+\zeta,               & c&=kY+\eta,\\
 a&=E+Y,                      & H&=4a+3,\\
 \Delta&=a^2+H=(a+2)^2-1,     & \ell&=EkY,\\
 b_0&=X+ac,                   & G&=1+\Delta c^2,\\
 R&=[(q-1)(q-F)+\alpha](q^2-1)+(M_C+qM_F)J,\\
 r&=k-R-1,                    & U&=Y^2,\\
 D_*&=Ub_0+H(U\gamma-r),       & M&=D_*^2-U^2G,\\
 I&=ic^2,                     & A_{\mathrm{aux}}&=\Delta(f^2-1),\\
 W&=I^2-A_{\mathrm{aux}},      & V&=f^2(cT-R)-c,\\
 y&=V+e,                      & N_a&=A_{\mathrm{aux}}(V^2-y^2)+y^2,\\
 L_t&=(K_0+w)C-F-z(q-1),\\
 A_0&=[q^3(\ell+\rho)+Y]^2-q^6[\ell(\ell+k)+1].
\end{aligned}
\end{equation}
The equation is
\begin{equation}\label{eq:sc-polynomial}
 \boxed{P=A_0^2+M^2+W^2+N_a^2+L_t^2-1=0.}
\end{equation}
Only fixed nonnegative powers occur.  The definition contains no division,
variable exponent, recurrence, bit test, or radical.  The literal evaluator
is \path{polynomial.json}; \path{build.py} reconstructs it.  Positivity is
part of the quantified domain in Theorem~\ref{thm:sc-main}.

\subsection{The inherited interface used in the proof}

For an integer $A\ge2$, define the Pell polynomials by
\[
 \chi_A(0)=1,\quad\psi_A(0)=0,\qquad
 \begin{aligned}
 \chi_A(j+1)&=A\chi_A(j)+(A^2-1)\psi_A(j),\\
 \psi_A(j+1)&=\chi_A(j)+A\psi_A(j).
 \end{aligned}
\]
These sequences appear only in the proof.  We use the standard Pell
classification, their strong divisibility, and the elementary recurrence
congruences.  The source proofs are preserved in
\path{upstream/PELL_LEMMAS.md} and
\path{upstream/ORDINARY_KERNEL.md}.

The particular compiler interface needed here is as follows.  Take fifteen
positive coordinates
\[
 (J,F,\alpha,z_{\rm old},f,h,i,T_{\rm old},s,w,
       \tau,\eta,\zeta,y_{\rm old},\gamma_{\rm old}),
\]
put $Y=sq^3$, use the shared abbreviations above, and set
\[
 D=b_0+H\gamma_{\rm old},\qquad
 V_{\rm old}=c(T_{\rm old}f-1)-Rf^2.
\]
The ordinary strong kernel consists of the six equations
\begin{equation}\label{eq:sc-source}
\begin{aligned}
 \tau^2-\ell(\ell+k)&=1, & D^2-\Delta c^2&=1,\\
 A_{\rm aux}(V_{\rm old}^2-y_{\rm old}^2)+y_{\rm old}^2&=1,
       & k-hE-R&=1,\\
 (K_0+w)C+(q-F)-z_{\rm old}(q-1)&=1,
       & (ic^2)^2&=\Delta(f^2-1).
\end{aligned}
\end{equation}
For an authentic complete tuple, a positive solution of this system gives
a finite blank K215 computation, hence a finite ZFC refutation.  Conversely
a finite refutation supplies canonical normalized compiler witnesses, which
give ordinary witnesses by
\begin{equation}\label{eq:sc-normalized}
 i=\Delta i_{\rm norm},\qquad
 f^2=1+\Delta i_{\rm norm}^2c^4.
\end{equation}
The canonical outer identities are
\begin{equation}\label{eq:sc-canonical}
 p=R=2n-1,\qquad R\equiv3\pmod4,\qquad
 X=wq=2^R,\quad Y=sq^3,\quad c=\psi_{a+2}(R).
\end{equation}
These statements are the ordinary compiler theorem in
\path{upstream/ORDINARY_KERNEL.md}, composed with the complete native
adapter in \path{native/PROOF.md}.  They specify the imported semantic
dependency precisely.  Sections below prove the new reduction to and from
this interface.

\subsection{A minimum block replaces an extra squared residual}

\begin{lemma}\label{lem:sc-minimum}
For every positive assignment to the fourteen coordinates,
\[
 P=0\quad\Longleftrightarrow\quad
 A_0=M=W=L_t=0\quad\hbox{and}\quad N_a=1.
\]
\end{lemma}
\begin{proof}
We have $a>0$, $\Delta>0$, $c\ge3$, and $I=ic^2\ge9$.
If $f=1$, then $A_{\rm aux}=0$ and $W=I^2>1$, which is incompatible
with $P=0$.  If $f\ge2$, then $A_{\rm aux}>1$.  In this case $N_a=0$
would imply
\[
 A_{\rm aux}V^2=(A_{\rm aux}-1)(V+e)^2.
\]
Neither $V$ nor $V+e$ could vanish, since simultaneous vanishing would
force $e=0$.  Thus $A_{\rm aux}(A_{\rm aux}-1)$ would be a rational
square.  A rational square root of an integer is integral, whereas
\[
 (A_{\rm aux}-1)^2
 <A_{\rm aux}(A_{\rm aux}-1)<A_{\rm aux}^2.
\]
This contradiction shows $N_a^2\ge1$.  The sum in
\eqref{eq:sc-polynomial} therefore forces the other four residuals to
vanish and $N_a=\pm1$.

With $W=0$ we have $A_{\rm aux}=I^2$, and
\[
 N_a=I^2V^2-(I^2-1)y^2.
\]
Modulo four this is $V^2$ if $I$ is odd and $y^2$ if $I$ is even.
It cannot be $-1$.  The converse follows by substitution.
\end{proof}

\subsection{Scale recovery and the adaptive modulus}

Suppose $P=0$.  From $A_0=0$, the rational number
\[
 \frac{q^3(\ell+\rho)+Y}{q^3}
\]
has square $\ell(\ell+k)+1\in\Z$.  It is integral, so $q^3\mid Y$.
Consequently
\begin{equation}\label{eq:sc-scale}
 s=Y/q^3>0,\qquad
 \tau=\ell+\rho+s>0,\qquad
 \tau^2-\ell(\ell+k)=1.
\end{equation}
In particular $Y\ge q^3$, $X\ge q$, and the authentic value of $\beta$
gives $q\ge16$.

Put $z_{\rm old}=z+1$.  The transport identity is
\[
 (K_0+w)C+(q-F)-z_{\rm old}(q-1)=1+L_t=1.
\]
It forces $C>0$: if $C\le0$, its left side is at most
\[
 q-F-z_{\rm old}(q-1)
   =1-F-(z_{\rm old}-1)(q-1)\le0.
\]
Hence $\alpha+C=q-F$ and $1\le F\le q-1$.  Set
$G_0=(q-1)(q-F)+\alpha=q\alpha+(q-1)C$.  Positivity gives
\[
 q\le G_0\le q^2-q-1,\qquad
 0<(M_C+qM_F)J<(q-1)(1+2q).
\]
For the upper bound on $G_0$, use $F\ge1$, $C\ge1$, and
$\alpha\le q-F-1$.  Since $J=(q-1)/(B-1)$, the mask bound follows from
the complete tuple's inequalities.  These estimates imply
\begin{equation}\label{eq:sc-bounds}
 3q+1<R<q^4-q^3<q^4-2,
 \qquad E\ge q^4>R+2,\quad a>E>R,
 \quad kY<c<k(Y+1).
\end{equation}

\begin{lemma}[Adaptive modulus]\label{lem:sc-modulus}
After the scale recovery \eqref{eq:sc-scale}, the polynomial $U=Y^2$
satisfies
\[
 U\mid4XY^2,\qquad
 \mu:=\frac{U}{\gcd(U,H)}\ge\frac{Y^2}{9}>R+2.
\]
Neither conclusion requires a value for the main Pell index or a
recovered index quotient.
\end{lemma}
\begin{proof}
The first divisibility holds identically.  Because
$H=4Y(X+1)+3$, we have $\gcd(Y,H)\mid3$.  Prime by prime,
\[
 \min(2v_t(Y),v_t(H))\le2\min(v_t(Y),v_t(H)),
\]
so $\gcd(Y^2,H)\mid9$.  Therefore
\[
 \mu\ge Y^2/9\ge q^6/9>q^4>R+2,
\]
where $q\ge16$ supplies the strict middle inequality.
\end{proof}

The distinction between size and degree is decisive.  Earlier
fourteen-variable quotient absorption used $U=4q^4$.  Here the directly
supplied coordinate $Y$ has degree one; integrality forces $Y\ge q^3$
only on a zero.  Thus $Y^2$ has smaller polynomial degree while meeting
the required arithmetic lower bound.

\subsection{The early bound precedes strong-rank recovery}

The equation $M=0$ makes $g=D_*/U$ a rational square root of the integer
$G$, hence an integer.  It is positive \emph{before} either Pell index
has been identified.  Indeed $r=k-R-1<k$, $a>3$, and $H<5a$, while
\[
 \frac{b_0}{H}>\frac{ac}{H}>\frac c5>\frac{kY}{5}
       >\frac{k}{Y^2}=\frac kU.
\]
The last inequality uses $Y^3>5$.  Hence
\[
 g=b_0+H(\gamma-r/U)>b_0-Hk/U>0,
 \qquad g^2-\Delta c^2=1.
\]
Reducing $D_*=Ug$ modulo $U$ gives
\begin{equation}\label{eq:sc-weakdiv}
 U\mid H(k-R-1),\qquad \mu\mid k-R-1.
\end{equation}
At this stage it is not legitimate to replace $\mu$ by $U$ in the second
divisibility.

Let $A=a+2$ and $P_0=2XY^2+1$.  The first norm is
$\tau^2-XY^2(XY^2+1)k^2=1$.  Its fundamental positive unit has second
coordinate two: coordinate one lies strictly between consecutive
squares, and coordinate two is attained by $2XY^2+1$.  Pell
classification, and then the main norm, give indices $n,p\ge1$ with
\begin{equation}\label{eq:sc-indices}
 \tau=\chi_{P_0}(n),\quad k=2\psi_{P_0}(n),
 \qquad g=\chi_A(p),\quad c=\psi_A(p).
\end{equation}
Since $P_0>A$ and $c>k$, strict growth gives $p>n$.
The recurrence congruence $\psi_{P_0}(n)\equiv n\pmod{P_0-1}$ yields
\begin{equation}\label{eq:sc-firstcong}
 k\equiv2n\pmod{4XY^2}.
\end{equation}
By Lemma~\ref{lem:sc-modulus} and \eqref{eq:sc-weakdiv},
$2n\equiv R+1\pmod\mu$.  Since $0<R+1<\mu$ and $n>0$, we can write
\begin{equation}\label{eq:sc-early}
 2n=R+1+t\mu\quad(t\ge0),\qquad n\ge(R+1)/2\ge25,
 \qquad p\ge26.
\end{equation}
This proves the required early bound without the deleted equation
$k-hE-R=1$.  It also gives
\begin{equation}\label{eq:sc-large}
 c>A\Delta^2,\qquad c>2p,\qquad c>2R.
\end{equation}
For the first inequality, use
$\psi_A(p)\ge(2A-1)^{p-1}\ge A^{25}>A(A^2-1)^2$.
The recurrence gives $c>2p$; finally $c\ge\psi_A(2)=2A>2R$.

\subsection{The ordinary rank and the auxiliary index}

We record the exact inherited rank statement, because a weaker
square-root observation would not suffice.

\begin{lemma}[Ordinary strong rank]\label{lem:sc-rank}
Let $A\ge2$, $\Delta=A^2-1$, and
$c=\psi_A(p)>A\Delta^2$.  If $i,f>0$ and
$(ic^2)^2=\Delta(f^2-1)$, then for an integer $m>0$,
\[
 f=\chi_A(m),\quad p\mid m,\quad c\mid m,\quad
 I=\Delta\psi_A(m),\quad c^2\mid(f^2-1).
\]
In the situation above, $m\ge c>2p$.
\end{lemma}
\begin{proof}
This is the ordinary strong-rank argument preserved in
\path{upstream/ORDINARY_KERNEL.md}, Section 4.  For completeness, its
essential rank step is recalled.  In the quadratic field, let $F_0$ be
the abscissa of the fundamental positive integer-coordinate Pell unit.
For some $d>0$ and $S=\psi_{F_0}(d)$,
\[
 A=\chi_{F_0}(d),\qquad
 \Delta=(F_0^2-1)S^2,\qquad
 \psi_{F_0}(dp)=Sc.
\]
Pell classification applied to the strong equation, using the squarefree
discriminant, gives $f=\chi_{F_0}(u)$ and
$I=(F_0^2-1)S\psi_{F_0}(u)$ for some $u>0$.
Put $b=\psi_{F_0}(\gcd(u,dp))$.
Strong divisibility and $c^2\mid I$ imply $b\ge c/\Delta$.
Indeed, if $g_1=\gcd(c,\psi_{F_0}(u))$, then
$c/g_1\mid(F_0^2-1)S$, while $b\ge g_1$ because
$c\mid\psi_{F_0}(dp)$.  Thus
$b\ge c/((F_0^2-1)S)\ge c/\Delta$.
If $\gcd(u,dp)$ were a proper divisor of $dp$, then
\[
 2b^2<\psi_{F_0}(2\gcd(u,dp))\le Sc<Ac,
\]
contradicting $c>A\Delta^2$.  Thus $dp\mid u$; write $u=dm$,
so $p\mid m$ and composition gives $I=\Delta\psi_A(m)$.

Write $m=pt$.  From $I=ic^2$ and
$\psi_A(pt)=c\psi_{\chi_A(p)}(t)$, reduction modulo $c$ gives
$c\mid\Delta t$.  Here the factor $\chi_A(p)$ is a unit modulo $c$.
Also
$\psi_A(p)\equiv pA^{p-1}\pmod\Delta$, whence
$g_0:=\gcd(c,\Delta)=\gcd(p,\Delta)$ divides $p$.
Since $c/g_0\mid t$, we obtain $c\mid pt=m$.
Finally
\[
 f^2-1=\Delta\psi_A(m)^2=ic^2\psi_A(m),
\]
which proves the required divisibility.  This proof includes the case
where $A$ is not the fundamental unit for the squarefree discriminant.
\end{proof}

Apply the lemma using $W=0$ and \eqref{eq:sc-large}.  Set
$T_{\rm old}=fT>0$ and $y_{\rm old}=|V+e|$.  Then
\begin{equation}\label{eq:sc-aux}
 V=c(T_{\rm old}f-1)-Rf^2,\qquad
 I^2V^2-(I^2-1)y_{\rm old}^2=1.
\end{equation}
Neither $V$ nor $y_{\rm old}$ vanishes; in particular $y_{\rm old}>0$.
Pell classification with parameter $I$ supplies an index $j>0$ such that
\[
 I|V|=\chi_I(j),\qquad y_{\rm old}=\psi_I(j).
\]
The index $j$ is odd, since even-index abscissas are $\pm1$ modulo $I$.
For $j=2v+1$, the integer polynomial $Q_v$ satisfies
\begin{equation}\label{eq:sc-Q}
 \chi_t(2v+1)=tQ_v(t^2),\qquad
 Q_v(0)=(-1)^v(2v+1),\qquad
 Q_v(1-A^2)=(-1)^v\psi_A(2v+1).
\end{equation}
One may verify these identities from
$Q_0(X)=1$, $Q_1(X)=4X-3$, and
$Q_{v+1}(X)=(4X-2)Q_v(X)-Q_{v-1}(X)$.

The explicit expression for $V$ gives
\[
 V\equiv-c\pmod f,\qquad V\equiv-R\pmod c;
\]
the second uses $c^2\mid f^2-1$, which has now been proved by rank
recovery.  Since $I^2\equiv1-A^2\pmod f$, identities
\eqref{eq:sc-Q} imply $\psi_A(j)\equiv\pm c\pmod f$.
Squaring and using the doubling identity gives
\[
 \chi_A(2j)\equiv\chi_A(2p)\pmod{\chi_A(m)}.
\]
The signed Pell step-down therefore yields $j\equiv\pm p\pmod m$.
To see the required uniqueness, reduce the abscissa index modulo $2m$
and reflect it into $[0,m]$ using
\[
 \chi_A(t+2m)\equiv-\chi_A(t),\qquad
 \chi_A(2m-t)\equiv-\chi_A(t)\pmod{\chi_A(m)}.
\]
Because $0<2p<m$, a negative representative cannot have the displayed
positive residue: for $r,2p\le m-1$,
$\chi_A(r)+\chi_A(2p)\le2\chi_A(m-1)<\chi_A(m)$.
The representative $r=m$ has residue zero, and strict growth then forces
the positive representative to be $2p$.  Thus
$2j\equiv\pm2p\pmod{2m}$, as asserted.

Since $c\mid m$, this gives $j\equiv\pm p\pmod c$.
On the other hand $c\mid I$ and the constant term in
\eqref{eq:sc-Q} give $j\equiv\pm R\pmod c$.
Therefore $R\equiv\pm p\pmod c$.  The strict bounds
$0<p,R<c/2$ in \eqref{eq:sc-large} force
\begin{equation}\label{eq:sc-pR}
 p=R.
\end{equation}
This argument uses neither the missing index quotient nor a restored
main-root congruence.

\subsection{Restore the full modulus and finish soundness}

We now have $n<p=R$.  Hence
$|2n-R-1|<R<\mu$.  Its divisibility by $\mu$ in
\eqref{eq:sc-early} forces $2n=R+1$.
Return to the stronger first-Pell congruence
\eqref{eq:sc-firstcong}: it shows that
\[
 r=k-R-1=k-2n
\]
is divisible by the \emph{full} $4XY^2$, and therefore by both $E$ and
$U$.  Moreover $r>0$, since $n\ge2$ and $\psi_{P_0}(n)>n$.
Thus
\begin{equation}\label{eq:sc-recover}
 h=r/E>0,\qquad \gamma_{\rm old}=\gamma-r/U\in\Z.
\end{equation}
The equality $D_*=Ug$ now restores
$g=b_0+H\gamma_{\rm old}$, with the full modulus $H$.
Since $n\ge2$, $k\ge4P_0>X$ and $c>k>X$.
Also $\Delta>(a+1)^2$ implies $g>(a+1)c$, so
\[
 g-b_0>c-X>0,\qquad \gamma_{\rm old}>0.
\]
All fifteen positive coordinates in \eqref{eq:sc-source} have been
reconstructed: use \eqref{eq:sc-scale}, $z_{\rm old}=z+1$,
$T_{\rm old}=fT$, $y_{\rm old}=|V+e|$, and
\eqref{eq:sc-recover}.  Their six source equations hold exactly.
The inherited compiler interface therefore gives a finite ZFC refutation.

The order of the argument is essential.  Initially only
$U\mid Hr$ and $\mu\mid r$ were available.  The bound on $\mu$ first
made the rank theorem applicable; rank recovery gave $p=R$; only then
did the first norm restore $U\mid r$ and the full main-root congruence.

\subsection{Canonical completeness and positive replacement witnesses}

Conversely suppose a finite ZFC refutation exists.  Choose canonical
normalized compiler witnesses as in
\eqref{eq:sc-normalized}--\eqref{eq:sc-canonical}, and replace the strong
coordinate by $i=\Delta i_{\rm norm}$.  Retain the outer coordinates.
The identity $R=2n-1$ and \eqref{eq:sc-firstcong} show that
$r=k-R-1>0$ is divisible by $4XY^2$, hence by $U=Y^2$.
Set $\gamma=\gamma_{\rm old}+r/U>0$.  Then $D_*=Ug$ and $M=0$.

The replacement for the first root also has a positive gap.  Since
$\ell=XY^2k>k$, the equation
$\tau^2=\ell(\ell+k)+1$ gives
\[
 \tau<\ell+k,
 \qquad \tau-\ell=\frac{\ell k+1}{\tau+\ell}>\frac k3.
\]
Canonical $n\ge2$ gives $k>XY$, and therefore
\[
 k/3>XY/3>Y/q^3=s.
\]
Thus $\rho=\tau-\ell-s$ is a positive integer.  Supplying $Y=sq^3$
directly makes $A_0=0$.

The transport shift is positive too.  Since $X=2^R$ and $R>3q$,
$w=2^R/q>q>F$.  As $C>0$, the old transport equation gives
\[
 (z_{\rm old}-1)(q-1)=(K_0+w)C-F>0.
\]
Hence $z=z_{\rm old}-1>0$ and $L_t=0$.

It remains to construct the auxiliary quotient and its positive gap;
an arbitrary old quotient need not be divisible by $f$.
With $I=\Delta i_{\rm norm}c^2$, choose fresh auxiliary coordinates by
\[
 V=\frac{\chi_I(R)}I,
 \qquad y_{\rm old}=\psi_I(R).
\]
The quotient is a positive integer because $R$ is odd.
Write $R=2v+1$; the canonical phase $R\equiv3\pmod4$ makes $v$ odd.
Equation \eqref{eq:sc-normalized} gives
$I^2=\Delta(f^2-1)$ and hence
$I^2\equiv1-A^2\pmod{f^2}$.  Applying the polynomial identities
\eqref{eq:sc-Q} modulo the \emph{square} modulus $f^2$ and modulo $c$
yields
\[
 V\equiv-c\pmod{f^2},\qquad V\equiv-R\pmod c.
\]
Also $\gcd(c,f)=1$ and $f^2\equiv1\pmod c$.
It follows that
\begin{equation}\label{eq:sc-Tconstruct}
 T=\frac{V+c+Rf^2}{cf^2}
\end{equation}
is a strictly positive integer, and its defining identity is exactly
$V=f^2(cT-R)-c$.

The Pell identity gives
$I^2V^2-(I^2-1)y_{\rm old}^2=1$.
Since $R\ge3$ and $I>1$,
$V\ge\chi_I(3)/I=4I^2-3>1$.  Thus
\[
 (I^2-1)(y_{\rm old}^2-V^2)=V^2-1>0,
\]
so $e=y_{\rm old}-V>0$.
Now $y=V+e=y_{\rm old}$, $W=0$, and $N_a=1$.
All fourteen supplied coordinates are positive; Lemma~\ref{lem:sc-minimum}
gives $P=0$.  This proves the converse and completes the semantic part of
Theorem~\ref{thm:sc-main}.

The use of the phase $R\equiv3\pmod4$ in this direction is justified
by the authentic native compiler.  In particular its low-mask condition
$(C-1)\mathbin{\mathrm{AND}}(M_CJ+1)=0$, together with
$q\equiv0\pmod4$, $J\equiv1\pmod4$, and $M_C\equiv2\pmod4$,
gives $C\equiv1\pmod4$.  Substituting in the packing expression for $R$
then gives $R\equiv C+2\equiv3\pmod4$.  This bit condition belongs to
the inherited proof, not the polynomial definition.

\subsection{Exact degree, evaluator cost, and comparison}

Assign degree one to each supplied coordinate and degree zero to fixed
integers.  Put
\[
 q_0=\beta J,\quad X_0=wq_0,\quad k_0=\eta+\zeta,
 \quad E_0=X_0Y,\quad c_0=k_0Y,\quad
 R_0=q_0^3(q_0-F).
\]
These are the relevant leading homogeneous forms.  The exact residual
degrees are
\[
 \deg A_0=12,\qquad \deg M=14,\qquad \deg W=10,
 \qquad \deg N_a=15,\qquad \deg L_t=2.
\]
The cancellations behind the upper bounds follow from exact polynomial
identities:
\begin{align*}
 A_0={}&q^6(2\ell\rho+\rho^2-\ell k-1)
           +2q^3Y(\ell+\rho)+Y^2,\\
 D_0={}&b_0+H\gamma,\\
 D_0^2-G={}&2ac(X+H\gamma)+(X+H\gamma)^2-Hc^2-1,\\
 M={}&U^2(D_0^2-G)-2UHrD_0+H^2r^2,\\
 N_a={}&V^2-(A_{\rm aux}-1)e(2V+e).
\end{align*}
For $M$, the first summand has degree at most $13$, and the other two
have degree $14$.  Their joint leading form is
\[
 \operatorname{in}_{14}M
       =8X_0^2Y^2R_0(k_0Y^3+2R_0),
\]
which is nonzero.  The leaders of $V$ and $A_{\rm aux}$ are
$-f^2R_0$ and $E_0^2f^2$, so
\[
 \operatorname{in}_{15}N_a=2E_0^2f^4eR_0.
\]
Only $N_a^2$ reaches degree $30$.  Therefore
\begin{equation}\label{eq:sc-leader}
 \operatorname{in}_{30}P
  =4\beta^{10}J^{10}w^4Y^4f^8e^2(\beta J-F)^2.
\end{equation}
For example the coefficient of $J^{12}w^4Y^4f^8e^2$ is
$4\beta^{12}\ne0$.  This proves exact degree $30$ with the actual
fixed coefficients, rather than only a syntactic degree bound.

The supplied evaluator uses $44+23+16=83$ binary operations.
Squaring and multiplication by a numeral are charged; reuse and fixed
leaves are free.  All fourteen witness ports, all four fixed ports, and
all $83$ gates are live.  This is an arithmetic-circuit count; it is not a
claim about bit complexity or the size of the fully expanded polynomial.

For orientation, the relevant pinned comparisons are
\begin{center}
\begin{tabular}{lrrr}
\hline
Construction & Positive variables & Degree & Operations\\
\hline
Package 55, fourteen-variable gap form & $14$ & $52$ & $79$\\
Package 58, free-ordinate form & $15$ & $30$ & $79$\\
This construction & $14$ & $30$ & $83$\\
Package 51, shortest evaluator & $15$ & $157$ & $66$\\
\hline
\end{tabular}
\end{center}
Against the previous fourteen-variable point, the degree decreases by
$22/52=42.3077\ldots\%$ for four additional gates.  Against the
fifteen-variable degree-$30$ point, it deletes one positive unknown for
four gates.  These are distinct degree, variable-count, and evaluator-cost
tradeoffs; the result makes no global minimality claim.

\subsection{Provenance and independently reproducible checks}

The comparison is pinned to ProveIt commit
\begin{center}
\path{beb0192cb8dc36f78c369c080f2adfbc1887a486}.
\end{center}
\path{SOURCE_PINS.json} records forty exact copied upstream files, their
repository paths and revisions, Git blob identities, and SHA256 hashes.
The earlier source revision required by the compact coefficient generator
is recorded separately in \path{native/SOURCE_PINS.json}.  The original
notices and license files accompany the native material.

The inherited ingredients are the compact native compiler (package 44),
ordinary strong-rank and signed auxiliary-index arguments, index-quotient
absorption (package 50), the auxiliary gap minimum (package 55), the
squared-modulus quotient (package 51), and integral-root scale recovery
(package 58).  The new ingredient is the choice $U=Y^2$, its early gcd and
size lemma, and its integration with those changes.  Its proof in
Lemma~\ref{lem:sc-modulus} does not assume a decoded power-of-two value
for $X$ or a previously recovered index quotient.

\path{verify.py} is independent of the circuit builder.  It checks source
and coefficient hashes, circuit wiring and liveness, exact local sparse
polynomial identities, independent scalar transcriptions, and entire
univariate polynomial evaluations of the circuit modulo three primes.
It also verifies the exact formal leading forms used in
\eqref{eq:sc-leader}.  Both ordinary and optimized Python executions pass.
The results are recorded in \path{CHECKS.json}.  The complete native
coefficient generator and independent auditor were also rerun in an
isolated copy; all four raw payload hashes agree with the delivered
integers, as recorded in \path{NATIVE_REBUILD.json}.

These finite checks authenticate the delivered implementation and the
explicit finite certificate data.  The universal existence equivalence
rests on the mathematical proof above and the stated compiler interface.
No finite check here decides whether a positive zero or a ZFC refutation
exists.


% Included section; theorem environments and \Z, \N, \Con, \Pos are supplied
% by the containing report.  All artifact paths below are relative to
% integer41/.  This file deliberately contains no preamble.
\section{A quartic in 41 unrestricted integer unknowns}

The main construction in this section is an explicitly specified polynomial
\(Q_{41}\in\Z[x_1,\ldots,x_{41}]\) of total degree four.  It reduces the
45 integer variables and 28 quadratic residuals of the inspected source
quartic to 41 integer variables and 26 residuals.  The first saving comes
from recovering the sign of the radix \emph{after} the Pell equations have
forced a power of two.  The other three come from using a larger scale to
exclude a negative common sign, so that two positivity representations can
be replaced by one representation of a positive product.

\begin{theorem}
For the four fixed integer coefficients specified below, ordinary
mathematics formalizable in ZFC proves
\[
 \exists\mathbf x\in\Z^{41}\;Q_{41}(\mathbf x)=0
 \quad\Longleftrightarrow\quad
 \text{there exists a finite ZFC refutation}.
\]
Consequently ZFC proves that the absence of integer zeros of \(Q_{41}\) is
equivalent to \(\Con(\mathrm{ZFC})\).  Exact elimination of polynomially
defined coordinates also gives an equivalent sextic in 36 integers and an
equivalent octic in 31 integers.
\end{theorem}

The statement uses the inherited native compiler, Pell lemmas, and
calibration of the literal blank K215 machine to syntactic ZFC
refutability.  Those interfaces are identified below.  This is not a new
Lean or Rocq kernel certificate and does not assert a global minimum.
All artifact paths in this section are relative to
\path{integer41/}.  The inspected ProveIt snapshot is
\path{beb0192cb8dc36f78c369c080f2adfbc1887a486}.

\subsection{The complete fixed coefficient tuple}

The coefficient ports are fixed numerals, not quantified variables.  Use
the \emph{whole} source34 tuple named
\path{mixed-parity-unrounded-1064921}, including its shifted high mask.
Put
\[
 d=1064921=121\cdot8801,\qquad
 \rho=2^{121},\qquad B=2^d,\qquad \beta=B-1.
\]
Here is a finite recipe that fixes the other three ports completely.  Let
\(x_{it}\) be the ordered west, north, east, and south labels of tile
\(i\), in the literal 943-entry list
\path{upstream/wang943_tiles.json}.  Let
\[
 (o_0,o_1,o_2,o_3)=(1020,1000,980,960),\qquad
 e_t(j)=o_t+2j,\qquad W_j=2^{10+2j},
\]
and define the finite set
\[
 \mathcal E=\{0,\ldots,942\}
 \cup\{e_t(j):0\leq t<4,\ 0\leq j<10\}
 \cup\{1063,1064,1095,1096,2195,3296\}.
\]
For \(0\leq i<943\), set
\[
 c_i=1+\sum_{t=0}^3\sum_{j=0}^9
       \operatorname{bit}_j(x_{it}+512)W_{10t+j}
       +\sum_{j=40}^{43}W_j.
\]
The remaining coefficients are
\(c_{e_t(j)}=W_{10t+j}\),
\((c_{1095},c_{1096},c_{2195},c_{3296})=(W_{40},W_{41},W_{42},W_{43})\),
and \(c_{1063}=c_{1064}=0\).  Finally put
\[
 \mu=1022+\sum_{j=0}^{43}W_j,\qquad
 Z_0=(2^{119}-1)\sum_{e=8792}^{8798}\rho^e
       +(2^{27}-1)\rho^{8799},
\]
\begin{align*}
 D_C&=\rho^{40}+\rho^{41}+\rho^{1100}+\rho^{2200}
       +\sum_{e\in\mathcal E}c_e
            (\rho^{4393-e}+\rho^{5495-e}),\\
 D_R&=\rho,\\
 M_C&=B-1-\sum_{\substack{e\in\mathcal E\\e\ne1}}\rho^e
             -2\rho^{1063},\\
 M_{F,n}&=4+\mu(\rho^{4393}+\rho^{5495})
       +\sum_{j=0}^9(\rho^{1000+2j}+\rho^{1021+2j})
       +\rho^{2195}+\rho^{3296}+Z_0,\\
 K_0&=D_C+BD_R,\qquad M_F=M_{F,n}+\beta.
\end{align*}
In particular, the port \(M_F\) is \emph{not} \(M_{F,n}\).  This source
binding includes the exact native geometry and completeness controller
proved in \path{upstream/COEFFICIENT_PROOF.md}; agreement of four coarse
inequalities alone would not establish the compiler interface.

The actual coefficients satisfy
\[
 \beta\geq17,\qquad 0<M_C<\beta,
 \qquad \beta<M_F<2\beta,\qquad K_0\geq0.
\]
They are supplied as gzip-compressed minimal unsigned little-endian byte
strings in \path{coefficients/beta.bin.gz}, \path{coefficients/K.bin.gz},
\path{coefficients/MC.bin.gz}, and
\path{coefficients/MF_source.bin.gz}, respectively.  Their bit lengths are
1064921, 1065043, 1064921, and 1064922; their total uncompressed size is
532479 bytes.  The builder reconstructs all four integers from the literal
tile list and matches every full raw SHA-256 digest in
\path{upstream/PAYLOAD_BINDING.json}.  The file
\path{SOURCE_PINS.json} records the exact source Git blobs.  No coefficient
depends on waiting for a machine to halt.

\subsection{All variables and all 26 quadratic residuals}

There are 32 arithmetic coordinates
\[
\begin{gathered}
 J,F,\alpha,z,f,h,r,v,w_4,\tau,\eta,\zeta,y,\gamma,o,j,\\
 N,E,P,c,a,X,D,C_2,L,F_2,V,g,I,B_{\mathrm{aux}},\Delta,u
\end{gathered}
\]
and nine further coordinates, grouped as
\(\boldsymbol\xi=(\xi_1,\xi_2,\xi_3)\),
\(\boldsymbol\upsilon=(\upsilon_1,\upsilon_2,\upsilon_3)\), and
\(\boldsymbol\omega=(\omega_1,\omega_2,\omega_3)\).
Every one of these 41 coordinates ranges over \(\Z\).
The names \(v,w_4\) are the JSON coordinates
\path{v_scale}, \path{w_scale}; the triples correspond, in order, to
\path{FA_pos1}, \path{FA_pos2}, \path{FA_pos3};
\path{FC_pos1}, \path{FC_pos2}, \path{FC_pos3}; and
\path{eta_zeta_pos1}, \path{eta_zeta_pos2}, \path{eta_zeta_pos3}.

Use the following \emph{affine abbreviations}, which introduce no
additional unknowns:
\[
 q=\beta J+1,\quad k=\eta+\zeta,\quad C=q-F-\alpha,
 \quad s=2v+1,\quad i=2r+1,\quad H=4a+3.
\]
The positive ternary form is
\[
 \Pos(x,y,z)=1+x^2+y^2+z(z+1).
\]
It is strictly positive on \(\Z^3\) and represents every positive
integer.  For surjectivity, apply the three-square theorem to \(4n-3\):
exactly one of the three squares is odd, so division by four gives the
displayed form after permuting the coordinates.

Define the residuals literally as follows:
\begin{align*}
 r_1&=L-q^2+1, &
 r_2&=N-q(L+1),\\
 r_3&=X-4qw_4, &
 r_4&=a-E-sN,\\
 r_5&=E(X+1)-Xa, &
 r_6&=c-\eta-k(a-E),\\
 r_7&=P-k-E(c-\eta), &
 r_8&=\tau^2-P(P-k)-1,\\
 r_9&=D-X-ac-\gamma H, &
 r_{10}&=\Delta-a^2-H,\\
 r_{11}&=C_2-c^2, &
 r_{12}&=D^2-\Delta C_2-1.
\end{align*}
The remaining rows are
\begin{align*}
 r_{13}&=k-1-hE-(q-F)(N-q)+CL-(M_C+qM_F)J,\\
 r_{14}&=(K_0+4w_4)C-z(q-1)-F,\\
 r_{15}&=F_2-f^2,\\
 r_{16}&=V+c-fo,\\
 r_{17}&=V-cj-hE+k-1,\\
 r_{18}&=g+V^2-y^2,\\
 r_{19}&=I-iC_2,\\
 r_{20}&=B_{\mathrm{aux}}-I^2,\\
 r_{21}&=B_{\mathrm{aux}}-\Delta F_2+\Delta,\\
 r_{22}&=y^2-B_{\mathrm{aux}}g-1,\\
 r_{23}&=X-2u^2,\\
 r_{24}&=F\alpha-\Pos(\boldsymbol\xi),\\
 r_{25}&=FC-\Pos(\boldsymbol\upsilon),\\
 r_{26}&=\eta\zeta-\Pos(\boldsymbol\omega).
\end{align*}
Here \(\Pos(\boldsymbol\xi)\) means
\(\Pos(\xi_1,\xi_2,\xi_3)\), and likewise for the other triples.
The required single equation is
\[
 \boxed{\quad Q_{41}=\sum_{\nu=1}^{26}r_\nu^2=0.\quad}
\]
Over integers this is exactly the simultaneous vanishing of all 26 rows.
After expanding the affine abbreviations, every row has degree at most
two.  Only \(r_8\) contains \(\tau\), through its term \(\tau^2\), so
\([\tau^4]Q_{41}=1\).  Thus the degree is exactly four for the actual
fixed coefficient tuple.  The machine-readable specification is
\path{quartic41.json}.

\subsection{The inherited Pell interfaces and the proof boundary}

For \(b\geq2\), define the integer Pell sequences by
\[
 \chi_b(n)+\psi_b(n)\sqrt{b^2-1}
       =(b+\sqrt{b^2-1})^n\qquad(n\geq0).
\]
We use the ordinary Pell classification and the following two inherited
rank results, with their size hypotheses stated explicitly.

\begin{lemma}[Ordinary strong rank and signed auxiliary step-down]
Let \(b\geq2\), \(\delta=b^2-1\),
\(c_0=\psi_b(p)>0\), and let \(i\) be a nonzero odd integer.
If
\[
 c_0>b\delta^2,\qquad (ic_0^2)^2=\delta(f^2-1),
\]
then \(|f|=\chi_b(m)\) for an integer \(m\geq1\) satisfying
\(p\mid m\), \(c_0\mid m\), and \(m\geq c_0\).
Suppose also that \(c_0>2p,2R\), \(R>0\), and, with
\(\varrho=|i|c_0^2\), there are integers \(V,y\) satisfying
\[
 \varrho^2V^2-(\varrho^2-1)y^2=1,\qquad
 V\equiv\pm c_0\pmod{|f|},\qquad V\equiv-R\pmod{c_0}.
\]
Then \(p=R\).
\end{lemma}

The ordinary strong-rank input is the relaxed-rank lemma cited in
\path{upstream/SIGNED_DOMAIN_29.md}, under ``Ordinary-strong systems and
product''.  Its use for either sign of the scale is recorded in
\path{upstream/ODD_SCALE_27_PROOF.md}, Section 2.  For clarity, the second
part follows from the signed auxiliary proof as follows.  Auxiliary Pell
classification gives an odd index \(\ell=2j_0+1\) and an odd Chebyshev
quotient \(Q_{j_0}\) with
\[
 Q_{j_0}(1-b^2)=(-1)^{j_0}\psi_b(\ell),\qquad
 Q_{j_0}(0)=(-1)^{j_0}\ell.
\]
Reducing modulo \(|f|\), squaring the first exposed congruence, and
applying the plus step-down lemma with \(0<2p<m\) gives
\(\ell\equiv\pm p\pmod m\).  Reducing modulo \(c_0\) gives
\(\ell\equiv\pm R\pmod{c_0}\).  As \(c_0\mid m\), this yields
\(R\equiv\pm p\pmod{c_0}\).  The strict bounds
\(0<R,p<c_0/2\) force equality.  In particular, the sign of an original
coordinate \(c=\pm c_0\) is immaterial here.  One must establish the
ordinary strong-rank conclusion before inferring
\(f^2\equiv1\pmod{c_0}\); simply dividing its norm by \(\delta\) is not
a valid substitute.

The proofs below also use the ratio estimates and main-root congruences
in \path{upstream/ODD_SCALE_27_PROOF.md}, Sections 3--4, and in
\path{upstream/SIGNED_DOMAIN_29.md}, under ``Signed gamma''.  We record
the substitutions and the inequalities that make those applications
valid; no finite search is being substituted for an unbounded Pell lemma.

\subsection{A sign-free packing bound and a positive global scale}

At any common zero set
\[
 \begin{gathered}
 Q=|q|,\quad Y=sq^3,\quad A=a+2,\quad
 G=q(q-F)-C,\\
 R=G(q^2-1)+(M_C+qM_F)J.
 \end{gathered}
\]
The first 13 rows give, by polynomial identities,
\[
 \begin{gathered}
 N=q^3,\quad E=XY,\quad a=(X+1)Y,\quad c=kY+\eta,\quad
 \Delta=A^2-1,\quad k=R+1+hE,\\
 \tau^2-(XY^2)(XY^2+1)k^2=1,\qquad
 D^2-\Delta c^2=1,\qquad D=X+ac+\gamma H.
 \end{gathered}
\]
For example, substituting \(a=E+Y\) into \(r_5=0\) gives
\(E=XY\), without division.  Substituting \(L=q^2-1\) and
\(N=q^3\) into \(r_{13}\) gives \(k-1-hE-R\).
Rows 15--22 give the further exact identities
\[
 (ic^2)^2=\Delta(f^2-1),\qquad
 (ic^2)^2(V^2-y^2)+y^2=1,\qquad
 V=of-c=jc-R.
\]

The two guards \(F\alpha>0\) and \(FC>0\) force
\(F,\alpha,C\) to have the same nonzero sign.  Their sum is \(q\),
so \(|q|\geq3\).  Since \(q=\beta J+1\), the case \(J=0\) is
impossible, and \(Q\geq\beta-1\geq16\).  This does not yet prove
\(q>0\).

If \(q=Q>0\), all three summands are positive and
\[
 Q+2\leq G\leq Q^2-Q-1,\qquad
 0<(M_C+qM_F)J<(2Q+1)(Q-1).
\]
If \(q=-Q<0\), put \(F_0=-F\), \(C_0=-C\).  Then
\[
 G=Q(Q-F_0)+C_0,\quad 2Q+1\leq G\leq Q^2-2,
 \quad 0<(M_C+qM_F)J<2Q(Q+1).
\]
These follow directly from the fixed mask bounds and
\(J=(q-1)/\beta\).  In the positive branch the resulting upper bound is
\(Q^4-Q^3\); in the negative branch it is
\(Q^4-Q^2+2Q+2\).  Consequently both branches satisfy
\[
 \boxed{\quad Q(Q^2-1)<R<Q^4-2,\qquad R>3Q+3.\quad}
\]
Neither \(\eta,\zeta\) nor transport has been used.

Since \(s\) is odd, \(Y\ne0\) and \(|Y|\geq Q^3\).  Row 23 gives
\(X\geq0\).  If \(X=0\), then \(a=Y\) and
\(D=ac+\gamma H\).  Reducing the main norm modulo \(H\), using
\(\Delta=a^2+H\), gives \(0\equiv1\pmod H\), impossible because
\(|H|=|4Y+3|>1\).  Thus \(X>0\).  The new scale row now gives
\[
 X=4qw_4\geq4Q,\qquad
 |E|=X|Y|\geq4Q^4>3R+1.
\]
The extra factor four costs no variable: \(w_4\) replaces the old
coordinate \(w\) through the affine substitution \(w=4w_4\).

\subsection{Recovering positivity of eta and zeta}

The last guard gives \(\eta\zeta>0\), so their signs agree and neither
is zero.  Suppose for contradiction that both are negative.  Write
\[
\begin{gathered}
 \eta_0=-\eta>0,\quad\zeta_0=-\zeta>0,\quad
 K=\eta_0+\zeta_0=-k>0,\\
 Z=|Y|,\quad
 b=|A|,\quad c_0=|c|,\quad P_0=2XY^2+1.
\end{gathered}
\]
For \(Y>0\) one has \(KY<c_0<K(Y+1)\); for \(Y<0\) one has
\(K(Z-1)<c_0<KZ\).  In either case \(c_0>K\), \(b\geq2\), and
\(P_0>b\).  The last inequality follows, for example, from
\[
 2XZ^2+1-\big((X+1)Z+2\big)>0
 \quad(X\geq64,\ Z\geq4096).
\]
The first and main Pell classifications therefore give
\[
 K=2\psi_{P_0}(n),\quad |\tau|=\chi_{P_0}(n),\qquad
 c_0=\psi_b(p),\quad |D|=\chi_b(p),\qquad n,p\geq1.
\]
Monotonicity implies \(p>n\): if \(p\leq n\), then
\(c_0\leq\psi_b(n)\leq\psi_{P_0}(n)=K/2\), a contradiction.
Because \(P_0\equiv1\pmod{|E|}\), the Pell recurrence gives
\(\psi_{P_0}(n)\equiv n\pmod{|E|}\).  The index row and \(k=-K\)
give
\[
 2n\equiv-(R+1)\pmod{|E|}.
\]
As \(2n>0\) and \(0<R+1<|E|\), this entails
\[
 2n\geq |E|-R-1>2R,\qquad\text{hence }p>n>R.
\]
In particular \(p\geq6\).  The elementary Pell growth bound
\(\psi_b(p)\geq(2b-1)^{p-1}\) yields
\[
 c_0>b(b^2-1)^2,\qquad c_0>2p,\qquad c_0>2R.
\]
For the first inequality it suffices to observe that
\((2b-1)^5>b^5>b(b^2-1)^2\) for \(b\geq2\); the second follows
from \(3^{p-1}>2p\) for \(p\geq6\).

Every hypothesis of the ordinary strong-rank and auxiliary lemma now
holds.  In particular \(i=2r+1\) is nonzero and odd, the two exposed
quotients give \(V\equiv-c\pmod{|f|}\) and
\(V\equiv-R\pmod{c_0}\), and \(c=\pm c_0\) is permitted.  The
lemma gives \(p=R\), contradicting \(p>n>R\).  Therefore
\[
 \eta>0,\qquad\zeta>0.
\]
This argument is the reason for the larger scale.  The weaker bound
\(|E|>R+2\) from the unscaled construction does not supply the
negative-index contradiction \(n>R\).

\subsection{Delayed radix sign recovery: the intermediate quartic44}

It remains to justify that the other signed coordinates cannot introduce
spurious roots.  For this purpose consider the intermediate system
\path{quartic44.json}: replace \(w_4\) by one coordinate \(w\), use
\(X=qw\), and replace the product guard on \(\eta\zeta\) by two
separate ternary guards \(\eta=\Pos(\mathbf e)\),
\(\zeta=\Pos(\mathbf t)\).  It has 44 variables and 27 residuals.
All other rows are unchanged, in particular \(X=2u^2\).  Starting with
a zero of the 41-variable system, put \(w=4w_4\) and supply ternary
representations of the now positive \(\eta,\zeta\).  This gives a zero
of the intermediate system exactly.

For a zero of that intermediate system the preceding sign-free packing
bounds remain valid.  The same zero-scale exclusion gives
\[
 X\geq Q\geq16,\qquad |Y|\geq Q^3,\qquad
 |E|\geq Q^4>R+2,\qquad \eta,\zeta>0.
\]
These are the early scale hypotheses required by the inherited signed
Pell argument.  In particular no assumption that \(q>0\) or that
\(w\) is a square is needed to establish them.

Here are the rank and sign steps explicitly.  If \(Y<0\), put
\(Z=-Y\), \(b=(X+1)Z-2\), \(c_0=-c=kZ-\eta\), and
\(P_0=2XZ^2+1\).  Then
\[
 k(Z-1)<c_0<kZ,\quad P_0>b,\quad c_0>k,\qquad
 k=2\psi_{P_0}(n),\quad c_0=\psi_b(p),\quad p>n.
\]
Now the index congruence has the positive sign:
\(2n=R+1+t_0|E|\), \(t_0\geq0\).  Thus
\(n\geq(R+1)/2\geq24\).  Pell growth gives
\(c_0>b(b^2-1)^2\) and \(c_0>2p\); also
\(b=|E|+Z-2>R\), so \(c_0>2R\).
The inherited rank lemma and the exposed \(V\)-congruences give
\(p=R\).  Since \(n<p=R\) and \(|E|>R+2\), there is no wrap:
\(t_0=0\), hence \(R=2n-1\).

Writing \(R=2r_0+1\), the negative-scale ratio estimate gives
\[
 Z>\frac{c_0}{k}\geq
 \frac12\left(\frac{(2b-1)^2}{2P_0}\right)^{r_0}
 >\frac{X^{r_0}}2>2^R.
\]
The strict middle inequality follows from
\[
 (2b-1)^2-2P_0X
 =4(2X+1)Z^2-20(X+1)Z+25-2X>0;
\]
for \(X\geq16\) this is positive at \(Z=4\), where it equals
\(46X+9\), and increases thereafter.  Here \(r_0\geq2\).
The two signs of the main root imply, respectively,
\[
 X\equiv-2^R\pmod H
 \quad\hbox{or}\quad X2^R\equiv1\pmod H.
\]
Both are impossible, because
\(0<X+2^R<X+Z<|H|\) and \(1<X2^R<XZ<|H|\).
Consequently \(Y>0\).

For \(Y>0\), the same rank argument applies with \(b=A\) and
\(kY<c<k(Y+1)\).  In particular the index row implies \(h\geq0\),
the Pell indices satisfy \(p>n\geq(R+1)/2\), and the growth conditions
of the ordinary strong-rank lemma again hold.  Thus \(p=R=2n-1\).
The positive-scale ratio estimate gives
\[
 a>\frac{X^{r_0+1}}3,\qquad 0<X2^R<a<H-1.
\]
If \(D=-\chi_A(R)\), the main-root recurrence would give
\(X2^R\equiv-1\pmod H\), contradicting that interval.
Hence \(D=\chi_A(R)\), and the recurrence gives
\(X\equiv2^R\pmod H\).  Both \(X\) and \(2^R\) lie strictly
between zero and \(H\), so
\[
 \boxed{X=2^R.}
\]
This is precisely where the sign of the radix can finally be recovered.

\begin{lemma}[Mersenne radix sign]
For \(d\geq3\) and \(t\geq0\),
\(-2^t\not\equiv1\pmod{2^d-1}\).
\end{lemma}
\begin{proof}
Reduce \(t\) modulo \(d\), obtaining \(0\leq t_0<d\).
The least positive residue of \(2^t\) is \(2^{t_0}\), whereas
\(1\leq2^{t_0}\leq2^{d-1}<2^d-2\).  It therefore cannot be the
residue \(-1\).
\end{proof}

Since \(q\mid X=2^R\), one has \(q=\pm2^t\) for some
\(0\leq t\leq R\).  The definition \(q=\beta J+1\) and the lemma
exclude the negative sign.  Thus \(q>0\), and common signs give
\(F,\alpha,C>0\).  Also \(J>0\), \(w=X/q>0\), and
\(s=Y/q^3>0\).  This argument delays the radix sign until it is justified;
it does not attempt to recover the omitted square condition on \(w\).

\subsection{Literal transport, compiler soundness, and completeness}

The literal transport identity is
\[
 (K_0+w)C=z(q-1)+F,\qquad w=4w_4
\]
in the 41-variable system.  Source57's accompanying prose contains a
coordinate typo: it writes \((z-1)(q-1)+F\) at this point, but that is
not the expansion of its literal row.  Accordingly the conclusion needed
here is \(z>0\).  Indeed \(R>3q\) implies \(2^R>q^2\), so
\(w=2^R/q>q>F\).  With \(C\geq1\) and \(K_0\geq0\), transport
gives \(z(q-1)>0\).  In the older unshifted positive-factor convention,
the correct coordinate is \(z_{\mathrm{source}}=z+1\); its transport
factor then equals one.

After power recovery, the compiler soundness used here is the stage
``Power and compiler recovery'' and the blank finite-window stage of
\path{upstream/FIXED_BLANK_67_PROOF.md}, transported to the full tuple
above by \path{upstream/COEFFICIENT_PROOF.md}.  Its inputs are the
recovered positive coordinates, the authentic Pell kernel,
\(q^3\mid Y\), the bounds \(3q+1<R<q^4\), and the literal transport
identity.  The half-binomial valuation gives the population threshold;
\(q=\beta J+1\) and the fact that \(q\) is a power of two imply
\(q=B^{n_{\mathrm{cell}}}\) for some positive integer
\(n_{\mathrm{cell}}\); inverse population and
transport give the disjoint native masks.  The tuple's complete native
geometry decodes a blank K215 accepting tableau.  Its inherited finitary
calibration gives a finite ZFC refutation.

The square condition on \(w\) and the condition that \(q\) be twice a
square are sufficient choices of canonical witnesses in source57, not
hypotheses of this post-power soundness stage.  A current solution need
not satisfy those two restrictions in its existing coordinates.  To
recover a source57 solution one may first decode the finite computation
and then choose fresh canonical witnesses.  This distinction is necessary
for the delayed-sign reduction.

Conversely, a finite ZFC refutation supplies a blank K215 halt and the
canonical witnesses of the full source34 adapter.  They satisfy
\[
 q=2^t\ (t\text{ odd}),\quad R\equiv3\pmod4,\quad R>t,
 \quad X=2^R,\quad w=2^{R-t},
\]
with positive \(F,\alpha,C,\eta,\zeta\) and positive odd \(s,i\).
The adapter proves these parity promises for its actual coefficients,
including its rank controller; an older power-of-five assumption is not
being transferred to the new exponent \(d\).
Because \(R-t\) is a positive even integer, it is at least two, and
\(4\mid w\).  Choose
\[
 w_4=2^{R-t-2},\qquad u=2^{(R-1)/2},\qquad
 v=(s-1)/2,\qquad r=(i-1)/2.
\]
Retain the other canonical arithmetic coordinates and supply ternary
representations of \(F\alpha\), \(FC\), and \(\eta\zeta\).
Every displayed residual vanishes.  This proves completeness and, with
the preceding soundness argument, the theorem.  No consistency premise
is needed in either the fixed coefficient recipe or the equivalence.

\subsection{Exact sextic and octic companions}

For a sextic, delete the five indicated definition rows and eliminate
their five coordinates everywhere by
\[
 N=q(L+1),\quad X=4qw_4,\quad C_2=c^2,
 \quad F_2=f^2,\quad B_{\mathrm{aux}}=I^2.
\]
The remaining 21 residuals in 36 integer variables have degree at most
three.  Their sum of squares, specified in \path{sextic36.json}, has
exact degree six, with
\[
 [I^4g^2]Q_{36}=1.
\]
The certificate comes from the retained residual \(y^2-I^2g-1\);
the complete sparse expansion checks the coefficient in the full sum.

For an octic, instead delete ten definition rows and eliminate
\[
 \begin{gathered}
 L=q^2-1,\quad N=q^3,\quad X=4qw_4,\quad
 C_2=c^2,\quad F_2=f^2,\quad B_{\mathrm{aux}}=I^2,\\
 g=y^2-V^2,\qquad P=E(c-\eta)+\eta+\zeta,\\
 D=4qw_4+ac+\gamma(4a+3),\qquad \Delta=a^2+4a+3.
 \end{gathered}
\]
All right-hand sides involve only retained variables.  The remaining 16
residuals in 31 integers have degree at most four.  Their sum of squares
in \path{octic31.json} has exact degree eight, with
\[
 [I^4V^4]Q_{31}=1,
\]
as verified in the full expansion; the contributing auxiliary residual
is \(I^2(V^2-y^2)+y^2-1\).

These are exact bijections of integer zero sets: projection forgets the
defined coordinates, and their polynomial definitions restore them
uniquely.  Each deleted row is a unit-coefficient definition and becomes
the zero polynomial after substitution.  For \(N\), the simultaneous
octic substitution uses \(q(L+1)=q^3\).  Thus neither companion needs
a new sign, Pell, or compiler argument.

\subsection{Reproduction, comparison, and limitations}

Run \path{build_verify44.py}, then \path{build_verify41.py}, then
\path{build_companions.py} with Python 3; all three also run under
\texttt{python3 -O}.  They use the standard library, reconstruct the fixed
integers, check their full hashes and source blob identities, check the
packing and literal transport identities, and verify exact substitution
and elimination.  They expand the sums of squares with symbolic fixed
ports, check that every listed coordinate is live, and certify the exact
degrees and displayed unit coefficients.  The resulting certificates
are in \path{certificates/}.  Finite diagnostics for the elementary bounds
and radix lemma supplement their proofs above; they do not certify the
universal Pell or compiler theorems by sampling.

The precise primary comparison is the inspected source57 quartic:
45 variables and 28 rows become 41 variables and 26 rows at degree four.
The sextic and octic improve the named historical 42-variable sextic and
36-variable octic comparisons to 36 and 31 variables, respectively.
These claims do not displace the collection's smaller-variable equations
of much higher degree, and no global minimality or established worldwide
novelty is asserted.  The contribution is the explicit delayed-sign
guard, the absolute-radix bootstrap, and the large-scale exclusion of a
negative common sign, composed with the inherited full coefficient and
compiler interfaces.  The mathematical proof has internal AI-assisted
checking; no new end-to-end proof-assistant certificate or external peer
review is claimed.  Attribution and the exact inherited licenses are
preserved in \path{NOTICE.md} and \path{licenses/}.


\section{A smaller certified packing in nine positive unknowns}
\label{sec:packing111}

The separate accompanying directory \texttt{nine-positive-111} supplies an
additional improvement to package 59's nine-variable equation. It retains
the complete original package-31 coefficient tuple and the same arithmetic
finalizer, but changes the finite isolation certificate. The new equation
has exact total degree \(1{,}781{,}068\) and a division-free evaluator with
\(470\) live binary operations: \(162\) additions, \(22\) subtractions,
and \(286\) multiplications. The previous point was
\(1{,}910{,}404\) and \(473\) operations. This is a verified finite
packing improvement, not a new universal variable bound.

The positive quadratic parent has 31 coordinates and 22 residuals. After
the inherited row changes \(r_{14}\leftarrow r_{14}+r_2\) and
\(r_{15}\leftarrow r_{15}+r_3\), negate row 22. All three row operations
are invertible. The new weights, in the literal source order, are
\[
\begin{gathered}
110,105,100,94,74,22,41,95,87,89,49,31,11,16,98,19,\\
97,4,20,8,2,3,35,47,111,30,90,63,78,103,26.
\end{gathered}
\]
The corresponding degree-separated kernels have degrees
\((1954,1617,1788)\), maximum coordinate weight \(W=111\), and
convolution height \(H=2010\). Their 89 records have exactly the intended
22 coefficient outputs: an exhaustive certificate checks all 22,544
source/expansion interactions. The highest quadratic kernel term is
\(+Z^{1788}\), arising from the negated last row.

Write \(M=2^{2134927}\), \(a=4096M=2^{2134939}\),
\(B=ab^2\), \(\ell=\sum_i B^{w_i}\), and \(U=\ell+g\).
For all positive \(b,g\), before decoding any digits, the new kernels obey
\[
D_2(B)>\tfrac12 B^{1788},\quad
|D_1(B)|<2MB^{1617},\quad |D_0(B)|<2MB^{1954},
\quad U\ge B^{111}.
\]
The two lower-layer ratios are bounded by \(2M/B^{282}\) and
\(2M/B^{56}\), each below \(1/8\). Thus the source word
\(G_s=D_2(B)U^2+D_1(B)U+D_0(B)\) is strictly positive on every outer
assignment. On the full soundness range \(g<4bB^{111}\),
\[
\frac{2G_s}{B^{2011}}<\frac{86Mb^2}{B-1}<\frac1{32}.
\]
This uses \(86\cdot32<4095\). The marked target maximum 1954 is 56 below
the convolution height, so the same sparse padding gives
\(S+T<B^{2123}\). After decoding, the coefficient bound is unchanged:
\(993Mb^2<B/4\). Consequently the inherited mask argument recovers exactly
the same positive source zero set.

The smaller code requires a new, exact valuation calibration:
\[
L=2bB^{2123},\quad R=L(S+T+1)-S-1,
\qquad N=2^{16012032}B^{1121}.
\]
The factor 2 in \(L\) and the absence of a factor \(b\) in \(N\) are
essential. When \(b=2^k\), \(B=2^v\), and \(v=2134939+2k\),
the low and high masks have combined population \(134v-31k-22\).
For \(u=\log_2L=1+k+2123v\),
\[
u+\operatorname{pop}(T)
=2257v-30k-21
=2(1121v+16012032)=2\log_2N.
\]
The binary two-block formula and Legendre's valuation identity therefore
give \(N^2\mid\binom{2R}{R}\) exactly when \(S\mathbin{\&}T=0\).
The \(k=0\) calibration is included and cannot create a spurious source zero.

Compose these registers with the unmodified free-coordinate Jones finalizer
of package 59. On the nine positive variables \(b,g,h,s,w,C,i,j,n\), its
output has the form
\[
\boxed{\ F_9=\bigl[\Theta^2-\Omega^2(a_0+64a_1)\bigr]^2
                   -256\Omega^4a_0a_1=0.\ }
\]
All registers and fixed integer leaves are defined in the accompanying proof
and literal circuit. In particular, \(j_0=j-1\) is retained, the complete
old tuple is used coherently, and no additional existential variable is hidden.
The positivity-first finalizer proof applies because its required inequalities
have just been reproved for this layout. Composing with the explicitly named
inherited compiler and calibration theorems yields the ordinary ZFC-provable
equivalence between absence of positive zeros and \(\operatorname{Con}(\mathrm{ZFC})\).

The exact leading degrees are \(r=\deg R=8491\) and
\(q=\deg N=2242\). Resolving the norm's cancellation gives
\[
\deg F_9=104r+400q+8W+316
=816H+1224W+5044=1{,}781{,}068.
\]
The diagonal leading coefficient is
\(2^{290}a^{890324}c^{400}>0\), with \(c=2^{16012032}\), so this
upper bound is attained for the actual coefficient tuple. The natural-domain
translation has the same degree and 477 operations. The supplied independent
checker verifies the finite certificate, fixed payloads, universal ring
identities, leading form, and circuit implementation. These finite checks do
not replace the inherited unbounded proofs, and the new literal equation has
not been certified end-to-end by Lean, Rocq, or an object-language ZFC kernel.


\section{Coefficient bindings and provenance}
\label{sec:coefficients}

The components use three complete tuples.  Sharing the symbol $\beta$, or
even its numerical value, never authorizes mixing the other three ports.
The two tuples with exponent $1{,}064{,}921$ have different transport and
mask data.  In every tuple $K_0$ denotes the stored $K$ port and $M_F$
denotes the already shifted \path{MF_source} port.  The shift by $\beta$
must not be applied again.

\begin{center}
\small
\begin{tabularx}{\textwidth}{@{}lXrr@{}}
\toprule
Component & Complete source tuple & $\log_2(\beta+1)$ & Raw bytes\\
\midrule
\path{positive14/} & Package 44, compact complement layout & $590{,}069$ & $295{,}388$\\
\path{integer41/} & Package 34, mixed parity unrounded layout & $1{,}064{,}921$ & $532{,}479$\\
\path{nine-positive-111/} & Package 31, original mixed layout & $1{,}064{,}921$ & $533{,}099$\\
\bottomrule
\end{tabularx}
\end{center}

The first four compressed payloads total only 9,445 bytes.  The short
compressed size is not used as a mathematical complexity claim; the raw
numerals remain large.  The new constructions do not reduce the coefficient
sizes.  The nine-variable component also contains twelve derived fixed
leaves for the finite source code, base, and threshold.

Every copied source is authenticated in a component
\path{SOURCE_PINS.json}.  These manifests record the actual path, source
commit, and Git blob identity, with SHA-256 where supplied.  A historical
\path{source_commit} field inside an unchanged inherited coefficient
manifest records that adapter's original starting point; it is distinct
from the review snapshot above.  All retained licenses and notices stay
with their sources.

\section{Verification, formalization, and integration}
\label{sec:verification}

\subsection{What the executable checks establish}

The root command is
\begin{center}\texttt{python3 verify\_all.py}\end{center}
from the extracted archive.  Python 3.10 or later and the standard library
are sufficient.  The driver checks the recorded baseline blobs, runs each
component generator and auditor in a temporary copy, verifies exact
regeneration of the literal polynomial data, and separately rebuilds the
native numerals.  It leaves the delivered source snapshots intact.

For $P_{14}$ the checker validates all fourteen live inputs, four fixed
leaves and 83 live gates; checks four local identities in an exact sparse
integer polynomial ring; resolves the relevant leading homogeneous forms;
and establishes the exact degree 30.  Modular transcription checks with the
actual coefficients are supplementary implementation diagnostics.  The gcd
and modulus theorem is proved for all integers in the text; its finite
sample checks are not its proof.

For the signed quartic, the checker reconstructs all four actual numerals
from the literal 943-tile list and checks the required inequalities.  It
checks the source packing and transport identities, every retained row
under the substitution $w=4w_{\mathrm{scale}}$, full symbolic expansion,
variable liveness, and the coefficient $[\tau^4]Q_{41}=1$.  Its 388-term
expansion treats the four fixed ports as separate symbols: this is not a
claim that exactly 388 collected monomials remain after all fixed integers
are substituted.  The companion eliminations are polynomial identities and
have their own exact degree certificates.

For $F_9$, the finite coefficient-extraction certificate exhausts all
22,544 source/expansion interactions and proves all 22 target identities.
It checks native and derived coefficient bytes, three exact ring identities,
the threshold calibration, circuit liveness and operation count, and the
leading form proving exact degree.  The search that found the positions
is not a proof dependency.  The fixed certificate and independent checker
are sufficient to reproduce the result.

The root driver can also run the new auditors with \path{--optimized} to
confirm that Python's \path{-O} mode does not disable acceptance checks.
Conditions use explicit errors rather than relying on \path{assert}.
The provided receipts state the exact scope of each check.

\subsection{What remains outside those finite checks}

The universal Pell classification, rank and congruence lemmas; the
native compiler's semantic soundness and completeness; the K215 endpoint;
and the inherited nine-variable finalizer are mathematical proof
dependencies.  Their relevant proof texts and source pins are included.
The new arguments check their hypotheses in the changed coordinates,
including the early index estimates and all sign branches.  No finite
evaluation can establish these universal existence equivalences on its own.

The active repository has a separately checked closed equation and checked
local arithmetic results.  Those files do not automatically certify these
new literal polynomials.  In particular, formalizing the native compiler
bridge and producing an object-language derivation remain distinct tasks
\cite{active}.  This delivery makes no claim to have completed either.
The written internal review is AI-assisted and is not external peer review.

\subsection{Useful next formalization targets}

The most localized new target is the adaptive-modulus theorem: prove the
integral-root scale recovery, $\gcd(Y^2,H)\mid9$, the early congruence bound,
and the later recovery of the full modulus $Y^2$.  These are exactly where
an unjustified cancellation could otherwise enter the argument.  Next,
formalize the signed packing bounds uniformly in $Q=|q|$, followed by the
Mersenne-radix sign lemma and the negative-common-sign contradiction.  The
remaining low-degree eliminations are exact polynomial substitutions.

For a repository integration, the self-contained directory can be placed
beside the existing numbered contributions.  It does not overwrite an older
package, replace the current formally certified equation, or require a
network connection to run its arithmetic checks.  The root README gives
the proposed measurements and the complete build instructions.  The
article builds by running \texttt{python3 build\_report.py}; the two LaTeX
passes use the included section sources.

\appendix
\section{Complete nine-variable arithmetic tail}
\label{app:tail}

For reference, the register definitions of the additional nine-variable
equation are reproduced here.  Its source code and kernels are defined by
the finite lists in \path{nine-positive-111/data/}; these supply the
$S,T,R,N,t$ used below.  The nine supplied coordinates are
$b,g,h,s,w,C,i,j,n\in\Z_{>0}$, and $j_0=j-1$.  All other symbols are
polynomial abbreviations, local to this appendix:
\begin{align*}
U&=N^2w,& Y&=N^2s,& M_0&=RY,\\
P&=2M_0^2U,& K&=R+1+h(P-1),& A&=M_0(U+1),\\
B_1&=2R+1,& W_1&=bw,& J&=4A-5,\\
Q&=3W_1C-2(W_1^2-1),&&&\\
D&=(A^2-1)C^2+1,& k_0&=iJC,& E&=k_0DC,\\
F&=(A^2-1)E^2+1,&&&\\
G&=A+F(CD+1-A),& H_0&=C+F(B_1+j_0C),&&\\
I&=(G^2-1)H_0^2+1,& a_0&=(P^2-1)K^2+1,& a_1&=DFI,\\
m&=K^2(4t-g)-16t(C-KY)^2,&&&\\
V_0&=PK,& T_{\mathrm{pair}}&=k_0(8D^2-4D-1),&V_1&=T_{\mathrm{pair}}GH_0,\\
\Omega&=Jm,& E_0&=J(n+W_1^2)+Q,&\Theta&=E_0-\Omega(V_0+V_1).
\end{align*}
The single equation is
\[
F_9=\bigl[\Theta^2-\Omega^2(a_0+64a_1)\bigr]^2
       -256\Omega^4a_0a_1=0.
\]
The exact polynomial is the literal division-free circuit in
\path{nine-positive-111/generated/split_111_positive.json}.  The companion
\path{split_111_natural.json} supplies the correct input shifts over
$\{0,1,2,\ldots\}$, with the same degree and 477 operations.  This is not a
nine-unrestricted-integer result.

\section{Archive entry points}

\begin{longtable}{@{}p{0.39\textwidth}p{0.56\textwidth}@{}}
\toprule
Path & Purpose\\
\midrule
\endhead
\path{README.md} & Results, domain distinctions, reproduction and integration.\\
\path{report/report.tex}, \path{report/report.pdf} & Main editable article and compiled paper.\\
\path{positive14/polynomial.json} & Literal 83-gate polynomial in fourteen positive variables.\\
\path{positive14/PROOF.md} & Complete new two-way reduction, including canonical witness construction.\\
\path{integer41/quartic41.json} & All 41 integer variables and 26 quadratic residuals.\\
\path{integer41/PROOF41.md}, \path{PROOF44.md} & Product-guard and delayed-sign proofs.\\
\path{integer41/sextic36.json}, \path{octic31.json} & Exact elimination companions.\\
\path{nine-positive-111/PROOF.md} & Full source packing, positivity, valuation and degree proof.\\
\path{nine-positive-111/data/} & Literal parent, positions, targets and homogeneous kernels.\\
\path{review/} & Pinned baseline statements and independent internal mathematical review.\\
\path{certificates/} & Root verification and reproduction receipts.\\
\path{verify_all.py} & One-command finite audit, with optional optimized-mode replay.\\
\path{MANIFEST.sha256} & Integrity ledger of the final archive.\\
\bottomrule
\end{longtable}

\begin{thebibliography}{9}
\raggedright
\bibitem{active}
Vladimir Reshetnikov and contributors,
\emph{A Diophantine equation for syntactic consistency of ZFC},
ProveIt active research README, pinned commit
\texttt{beb0192cb8dc36f78c369c080f2adfbc1887a486}.
\url{https://github.com/VladimirReshetnikov/ProveIt/blob/beb0192cb8dc36f78c369c080f2adfbc1887a486/Computability/DiophantineConsistency/README.md}.

\bibitem{source55}
ProveIt contribution 55, \emph{Congruence norms}, fourteen-variable proof.
\url{https://github.com/VladimirReshetnikov/ProveIt/blob/beb0192cb8dc36f78c369c080f2adfbc1887a486/Computability/DiophantineConsistency/con-zfc-diophantine-forms/55-congruence-norms/fourteen/PROOF.md}.

\bibitem{source57}
ProveIt contribution 57, \emph{Separated degrees}, literal 45-integer quartic
and proof.
\url{https://github.com/VladimirReshetnikov/ProveIt/tree/beb0192cb8dc36f78c369c080f2adfbc1887a486/Computability/DiophantineConsistency/con-zfc-diophantine-forms/57-separated-degrees/quartic}.

\bibitem{source58}
ProveIt contribution 58, \emph{Integral-root carriers for Con(ZFC)},
paper, native coefficient binding and verification files.
\url{https://github.com/VladimirReshetnikov/ProveIt/tree/beb0192cb8dc36f78c369c080f2adfbc1887a486/Computability/DiophantineConsistency/con-zfc-diophantine-forms/58-integral-root}.

\bibitem{source59}
ProveIt contribution 59, \emph{Degree-separated encodings}, including the
nine-positive free-coordinate adjunct.
\url{https://github.com/VladimirReshetnikov/ProveIt/tree/beb0192cb8dc36f78c369c080f2adfbc1887a486/Computability/DiophantineConsistency/con-zfc-diophantine-forms/59-degree-separated}.

\bibitem{coefficients}
ProveIt contributions 31, 34 and 44, native compiler coefficient adapters
and retained Pell proofs.  Exact source locations, Git blob identities,
payload digests, and copies of the applicable proofs are supplied in each
component's \path{SOURCE_PINS.json}, \path{upstream/}, and native binding.

\bibitem{jones}
James P. Jones, \emph{Universal Diophantine equation},
Journal of Symbolic Logic \textbf{47} (1982), 549--571.
\url{https://doi.org/10.2307/2273588}.
The nine-variable component uses the specific inherited finalizer and its
retained proof, rather than an unspecified application of this reference.
\end{thebibliography}
\end{document}
